\icode{N18-EME-SI-L2-261121-001}
\nceeID[6]{2018全国\Rmnum{1}}{16}
如图，三个固定的带电小球\(a\)、\(b\)和\(c\)，相互间的距离分别为\(ab = \)\qty{5}{cm}，\(bc =\)\qty{3}{cm}，\(ca =\)\qty{4}{cm}。小球\(c\)所受库仑力的合力的方向平衡于\(a\)、\(b\)的连线。设小球\(a\)、\(b\)所带电荷量的比值的绝对值为\(k\)，则\paren[]

\begin{tikzpicture}[>=stealth,font=\scriptsize]
	\draw[dashed,cyan] (37:3.2)coordinate(c)--(0,0)coordinate(a)--(4,0)coordinate(b)--cycle;
	\foreach \c/\n in {a/below left,b/below right,c/above}
		\fill (\c) node[\n]{\(\c\)} circle(.04);
\end{tikzpicture}
\begin{choices}
	\item \(a\)、\(b\)的电荷同号，\(k = \frac{16}{9}\)
	\item \(a\)、\(b\)的电荷异号，\(k = \frac{16}{9}\)
	\item \(a\)、\(b\)的电荷同号，\(k = \frac{64}{27}\)
	\item \(a\)、\(b\)的电荷异号，\(k = \frac{64}{27}\)
\end{choices}
\begin{answer}
	D\quad 提示：
\end{answer}